\chapter{Dataset and Annotation Methodology}\label{dataset-and-annotation-methodology}

\section{Main Flow Dataset}\label{main-flow-dataset}

The main evaluation data consists of 13 web interaction flows derived from the Mind2Web dataset. Mind2Web was originally introduced as a benchmark for generalist web agents and contains real website tasks and interaction trajectories \autocite{deng2023}. In this study, ordered screenshots from selected trajectories serve as observations for independently represented UI requirements; the original action-prediction task is not evaluated.

The selected flows cover theme-park purchases, product lookup, careers
navigation, dining information, cinema support forms, cruise search, book
search, and business listings. Each flow contains several meaningful states. A
cart flow may progress from product selection through quantity changes and
add-ons to the summary and checkout controls. Search flows contain inputs and
later results. These sequences support requirements that a single screenshot
cannot resolve.

The benchmark contains 258 verification items across the 13 flows: 173 reviewed
source requirements and 85 reviewed contrastive items. Their annotations
contain 541 claims. All items completed a prediction-independent primary-author
review workflow.

The label distribution is imbalanced: 172 items are \texttt{FULFILLED}, 45 are \texttt{PARTIALLY\_\allowbreak FULFILLED}, 33 are \texttt{ABSTAIN}, and 8 are \texttt{NOT\_FULFILLED}. Approximately two thirds of the items are therefore positive. A classifier that favors \texttt{FULFILLED} can obtain a deceptively strong accuracy while failing on the three less frequent labels. This imbalance motivates macro-F1 and per-class reporting.

UI evaluability is annotated separately from fulfillment. The 258 items include
192 \texttt{UI\_VERIFIABLE}, 62 \texttt{PARTIALLY\_\allowbreak UI\_\allowbreak VERIFIABLE}, and 4
\texttt{NOT\_UI\_VERIFIABLE} cases. Evaluability describes whether screenshots can
resolve the requirement in principle; the final label describes what the
recorded flow supports. An incomplete flow can leave a UI-verifiable
requirement undecidable, while a compound requirement with visible and hidden
parts can be only partially UI-verifiable.

\section{From Candidates to Verification Gold}\label{from-candidates-to-verification-gold}

Because requirements are derived from UI flows, benchmark construction uses
staged artifacts and prediction-independent human review to separate candidate
generation from reference annotation.

Harvesting produces a broad set of requirement hypotheses, including redundant,
ambiguous, hidden, and incompletely supported properties. Filtering and
rewriting convert selected hypotheses into candidates. A reviewer then rejects,
edits, or promotes each candidate; candidate status carries no correctness
claim.

The 13 Mind2Web flows currently contain 100 committed candidate requirements and 173 reviewed source requirements. These counts refer only to the Mind2Web-derived directories; combining them with PURE candidates produces a different total and must be avoided in reporting. Review promotes a text into the source requirement set, but the verification benchmark adds further information: UI evaluability, a requirement-level label, atomic claims, claim statuses, evidence steps, rationale, and uncertainty reasons.

Contrastive requirements are used to increase difficulty and improve label coverage. A contrastive item may strengthen a visible requirement with a completeness condition, add a hidden persistence obligation, require a comparison that the flow does not show, or request a control that is visibly absent. The automatically proposed contrast and intended label are treated as suggestions. They become benchmark reference data only after review against the requirement wording and the screenshots.

The resulting set contains more than straightforward positive descriptions, but
its construction introduces dependencies. Source requirements and verification
evidence originate from the same flows, and contrastive generation may
emphasize particular hard cases. The primary-author reference review was
performed independently of the evaluated model predictions. Source and
contrastive requirements are reported separately where their construction is
relevant, and flow-level resampling preserves the benchmark's clustered
structure.

\section{Annotation Schema}\label{annotation-schema}

Each verification item contains a requirement-level decision and claim-level detail. The four requirement labels are interpreted as follows:

\begin{itemize}
\tightlist
\item
  \texttt{FULFILLED}: all central observable claims are visibly supported, no central claim is contradicted, and at least one evidence item is present.
\item
  \texttt{PARTIALLY\_\allowbreak FULFILLED}: at least one important claim is supported while another important claim is missing, hidden, or ambiguous, with no central visible contradiction that would justify a negative label.
\item
  \texttt{NOT\_FULFILLED}: a central observable claim is contradicted by visible evidence.
\item
  \texttt{ABSTAIN}: the screenshots are insufficient for a reliable positive or negative decision, or the requirement primarily concerns non-visible properties.
\end{itemize}

Claims use statuses including \texttt{SUPPORTED}, \texttt{MISSING}, \texttt{CONTRADICTED}, \texttt{HIDDEN}, \texttt{AMBIGUOUS}, and \texttt{OUT\_OF\_SCOPE}. They are marked as core or supporting obligations and linked to screenshot steps. The claim layer is intended to expose why a compound requirement receives a partial or abstaining decision. It also creates a distinct evaluation problem: predicted claims may not have a one-to-one textual correspondence with gold claims and must be matched before their statuses are compared.

The primary evidence unit for the label experiments is a screenshot step.
Evidence-step annotation is stable across the benchmark and directly represents
the ordered-flow setting. A step annotation identifies which screen contains
the observation, while a textual note describes the relevant visible content.
Region-level grounding is evaluated in a separate evidence-localization
analysis because a correct box is a different construct from a correct
requirement label or screenshot-step trace.

\section{Running Example: Public Access to Amtrak Dining Information}\label{running-example-public-access-to-amtrak-dining-information}

The following requirement illustrates a clean positive case:

\begin{quote}
The system shall make onboard dining information and café menu resources discoverable through public site navigation without requiring the user to sign in.
\end{quote}

The requirement contains two observable claims: public navigation exposes
onboard dining and café-menu resources, and reaching them requires no sign-in.
Step 1 establishes the public site context. Step 4 shows the Onboard Dining
page and its route to Café content; step 5 shows the Café page and menu
resources. No sign-in wall appears along this path. Steps 1, 4, and 5 support
both claims.

The same screenshots do not resolve a stronger requirement that route-specific
menus appear only when applicable because no route context is visible. Account
ownership checks would also introduce hidden access-control behavior. Such
changes in wording can move the correct label from \texttt{FULFILLED} to
\texttt{PARTIALLY\_\allowbreak FULFILLED} or \texttt{ABSTAIN}; the written scope determines the decision.

\section{PURE as Exploratory External Material}\label{pure-as-exploratory-external-material}

PURE is a corpus of public requirements documents collected from heterogeneous sources and formats \autocite{ferrari2017}. Its documents are useful because their requirements are not generated from the screenshot trajectories used in the main benchmark. They are often longer, more formal, and dependent on headings, surrounding paragraphs, figures, or system context.

The pipeline extracts and contextualizes selected PURE requirements and
associates them with UI images from the documents. The Split/Merge subset
contains 31 accepted verification items and 78 claims, while Mashboot contains
11 accepted items. Every accepted item received manual review independently of
the evaluated model predictions. Candidate-generation provenance is stored
separately from the human acceptance decision and reference annotation.

PURE can support qualitative discussion about context dependence, compound
requirements, and UI verifiability. The analysis treats the associated UI
images as document artifacts and reports document-to-UI consistency separately
from implementation conformance over recorded execution flows.

\section{Review Workflow and Quality Controls}\label{review-workflow-and-quality-controls}

Review occurs at the candidate and verification stages. Candidate review checks
whether the text forms a meaningful, self-contained requirement for the
intended application context. Verification review checks UI evaluability,
label, claims, evidence steps, uncertainty reasons, and rationale against the
ordered flow.

The workbench validates structural constraints. A fulfilled item requires
evidence and may not contain a contradicted core claim. A negative item requires
visible counter-evidence. A partial item requires supported and unresolved
material content. An abstention requires an insufficiency reason. These checks
identify internally inconsistent records, but they cannot decide whether the
reviewer's interpretation of the screenshots is correct.

All 258 Mind2Web items and all accepted PURE items received manual
primary-author review independently of the evaluated model predictions. The
repository records review status and annotation provenance so that
candidate-generation assistance remains distinct from human reference
decisions. The completed primary-author decisions form the reference set used
in the reported evaluation.

UI-evaluability and region-grounding reviews use their own targeted samples.
Separate review tasks preserve differences between fulfillment, UI
evaluability, screenshot selection, and region sufficiency. Agreement at one
level does not establish agreement at another.

\section{Benchmark Characteristics}\label{benchmark-characteristics}

The benchmark is organized by flow rather than as an unstructured pool. Items
within one flow share application context and screenshot evidence, which
creates statistical dependence. Flow identifiers are preserved in every
prediction and metric artifact so that evaluation can resample or diagnose at
the flow level.

The 541 reviewed claims exceed the number of requirements because compound
items contain multiple obligations. The distribution of claims is not uniform:
some requirements are atomic, while others describe an interaction and its
result, several visible fields, or a visible property combined with a hidden
outcome. This variation makes raw-versus-decomposed comparison meaningful but
also complicates claim matching.

Reference evidence may contain one step or several steps. Multiple steps are
common where a requirement concerns navigation, persistence across a short
interaction, or the difference between an input state and a result state. The
reference set is not assumed to enumerate every semantically valid alternative.
Automated overlap metrics are therefore complemented by qualitative review.

Every frozen run manifest records the expected requirement count and hashes of
the gold files. Evaluation rejects or reports missing predictions rather than
silently treating them as abstentions. This guard is necessary because several
historical repository metrics were produced after the gold set had grown and
therefore used incompatible denominators.

\section{Data Governance}\label{data-governance}

The 13 Mind2Web flows belong to the separately distributed \texttt{test\_task} split.
Original screenshots, HTML, traces, videos, and unzipped test records remain
local. The repository can recreate the working flow set from identifiers after
the user obtains Mind2Web through the official source, but it does not mirror
the underlying test archive.

PURE source documents are likewise obtained from the original Zenodo record.
Because the curators note uncertainty about rights in the underlying documents,
the public artifact excludes PDFs, XML, figures, and substantial source text.
The full working repository is maintained under controlled access, while the
release candidate is limited to code, configurations, citations, and aggregate
results. Section 7.7 discusses the implications for reproducibility.
